\documentclass[11pt,a4paper]{article}

\usepackage[T1]{fontenc}
\usepackage[utf8]{inputenc}
\usepackage[magyar]{babel}
\usepackage{lmodern}
\usepackage[margin=2.4cm]{geometry}
\usepackage{microtype}
\usepackage{booktabs}
\usepackage{enumitem}
\usepackage{xcolor}
\usepackage{titlesec}
\usepackage[hidelinks]{hyperref}

\definecolor{accent}{HTML}{1F4E79}
\titleformat{\section}{\large\bfseries\color{accent}}{\thesection.}{0.6em}{}
\titleformat{\subsection}{\normalsize\bfseries}{\thesubsection}{0.6em}{}
\setlist{itemsep=2pt,topsep=3pt}

\title{\vspace{-1.2cm}\color{accent}\Large Heti összefoglaló --- HSiKAN, megfogás és struktúra}
\author{Hajdu Csaba dr. \\ \small Kati részére}
\date{2026.\ június 26.}

\begin{document}
\maketitle
\vspace{-0.6cm}

\noindent\textit{A heti munka a megfogás (grasping) / strukturált tanulás vonalon.
Megkülönböztetve, ami \textbf{mért}, ami \textbf{következtetett}, és ami még \textbf{hipotézis}.}

\section{A fő eredmény --- az ,,HSiKAN veszít az MLP ellen'' kérdés tisztázva}

A visszatérő aggály az volt, hogy a HSiKAN a robotikai RL-feladatokon legjobb esetben is csak
\emph{döntetlenre} hozza ki a sima MLP-vel szemben. Ez a hét lezárta a kérdést: \textbf{nem hiba}
--- sem a bekötésben, sem a backbone-ban.

\begin{itemize}
  \item \textbf{Bekötési audit (mért):} cartpole / arm6dof / négylábú / pick-place / FANUC --- mind
        tiszta. Nincs megfigyelés$\leftrightarrow$csúcs eltolódás, nincs elfajult sor, ép előreterjesztés.
        A cartpole HSiKAN megtanulja: 200/200.
  \item \textbf{Strukturális próba (mért):} felügyelt, RL-mentes, paraméter-illesztett HSiKAN vs.\ MLP.
        A backbone egészséges. \emph{Meglepetés:} a HSiKAN valódi előnye \textbf{elsősorban a poolingból}
        jön (csúcsonkénti aktiváció + átlag-pooling = Deep-Sets előfeltevés): egy szeparálható
        összegző célon akár \textbf{52$\times$}. Az előjeles 2-lépéses (signed 2-hop) előny másodlagos
        (kb.\ 1,5--3$\times$, adatigényes). Mindkét előny \emph{nő} az adattal $\Rightarrow$ reprezentációs,
        nem pusztán mintahatékonysági nyereség.
  \item \textbf{Újraértelmezés (következtetett):} a robotikai döntetlen egy
        \textbf{előfeltevés$\leftrightarrow$célfüggvény eltérés}, nem alkalmatlanság. A jelenlegi
        jutalmak mellett a robot érték-/policy-függvénye nem poololt/előjeles-strukturális alak, így a
        HSiKAN-nek nincs mit kamatoztatnia. \emph{Következtetés:} nem elvetni a HSiKAN-t, hanem olyan
        struktúrát adni neki, amit ki tud használni.
\end{itemize}

\section{Valódi jutalom-hiba a közös (kétkarú) jelenetben --- megtalálva és javítva}

Ez közvetlenül a közös megfogási feladaton van. A sűrű közelítési jutalom a \textbf{legközelebbi
kartestet} mérte (könyök, vagy rosszabb esetben a \emph{mozdíthatatlan} alaptest), nem az
\textbf{ujjhegyet} --- az emittált kar egyáltalán nem szállított ujjhegy-pontot (site). Egy seednél ez
\emph{konstans} értéket mért $\Rightarrow$ \textbf{nulla gradiens} az ujjhegy érme felé vezetéséhez.
Pontosan a megfigyelt tünet: ,,a robot pozícióját jutalmazzuk, nem a kar végét''.

\medskip
\noindent\textbf{Javítás:} tömeg nélküli \texttt{tip\_\{side\}} site-ok injektálása; a jutalom egy
ujjhegy-domináns keverék (0,75$\cdot$ujjhegy $+$ 0,25$\cdot$könyök). A fordított dinamika bitre azonos.
Egy változtatás három dolgot javított: a jutalom-metrikát, a BC-demonstrátort és a kar-hossz feloldást.
26 teszt zöld, köztük egy regresszió, amely a \emph{régi} metrikára megbukik.

\begin{table}[h]
\centering
\small
\begin{tabular}{cccc}
\toprule
seed & érme & régi (legközelebbi test) & valódi ujjhegy \\
\midrule
2 & (0,13;\,0,07) & 0,091 m (könyök)        & 0,220 m \\
4 & (0,19;\,0,07) & 0,076 m (könyök)        & 0,216 m \\
5 & ($-$0,04;\,0,06) & 0,230 m (\textbf{rögzített} alap) & 0,310 m \\
\bottomrule
\end{tabular}
\caption{A régi metrika ,,közelinek'' jelezte a kart (0,076--0,091 m), miközben az ujjhegy 0,22 m-re
volt --- a kar \emph{behajlítását} jutalmazta, nem a nyúlást.}
\end{table}

\section{További elkészült munka a héten (tesztelve, nem-core, CPU)}

\begin{itemize}
  \item \textbf{HTL-robusztusság mint jutalom (PoC):} egy \texttt{.htl} formula $\to$ egyszerre sűrű
        jutalom és monitor-ítélet, a meglévő HTL-kiértékelőt újrahasználva. Ez a híd az
        elszámoltathatósági réteg és a jutalom-formálás között.
  \item \textbf{RL-algoritmusvonal lezárva:} DDPG / TD3 / SAC immár a kódbázisban (a DDPG kb.\
        250$\times$ mintahatékonyság a PPO-hoz képest az off-policy vonalon).
\end{itemize}

\section{Javasolt következő lépés}

A kérdést végleg lezáró döntő kísérlet a \textbf{kiolvasó (readout) ablação}: a HSiKAN átlag-poolingját
--- ami a szeparálható összegeknél segít, de a vezérléshez szükséges \emph{ízületközi koordinációt}
összeomlaszthatja --- lecserélni konkatenáló / figyelmi (attention) kiolvasásra, egyetlen robotfeladaton.
Ha a döntetlen átfordul, a szűk keresztmetszet a pooling volt, nem a strukturális érvelés. Ezt egy
\textbf{strukturális robot-jutalommal} párosítva (feladatgráf-hiperélek + HTL-prediká\-tumok), hogy a
strukturális előfeltevésnek végre legyen olyan célfüggvény, amely jutalmazza.

\section*{Őszinteségi megjegyzés}

Két háttérfutás (HTL teljes A/B; a determinizmus-javítás utáni próba-söprés) újraindításkor leállt,
ezeket újra kell futtatni. A javított jutalmú közös szállítási arányt még nem mértük újra a 0,208-as
alaphoz képest. Tehát az ujjhegy-javítás egység- és dinamika-szinten igazolt, de a végponttól végpontig
mért javulás még függőben.

\end{document}
